

Traditional autoregressive language models (ARLMs) have shown remarkable performance (\cite{llama3herd2024}).
Diffusion language models (DLMs) have arisen as a new class of transformers. Instead of generating the output one
token at a time, from left to right, DLMs perform similar to their image generation counterparts, starting with noise
 and refining its predictions in parallel. Several new DLMs have reached similar performance at faster response times
  (\cite{nie2025large, google2025gemini, ye2025dream7b}. They also have built in error correction, as well as
  bidirectional context, which can enable more coherent outputs. Previous works have proposed ARLM watermarks,
  which bias the probability distribution of the model's predictions. The result is a signal imperceptible to humans
   but can be traced. ARLM watermarks have already been adopted in consumer models, and are advocated for with
   regulation. However, the unique architecture of DLMs has led to the development of two kinds of watermarking methods adapted for them.
   One scheme uses context hashing in expectation to pseudorandomly modify the token logit distribution \cite{gloaguen2025dlm}. The other method
   uses a keyed rule to reorder which positions are unmasked during diffusion generation \cite{hong2026dgmark}. We unify these two techniques and
   contribute a joint watermark that perturbs both token predictions at indices and the confidence metric used to select unmasking order, adding a signal
   inperceptible to the human eye but detectable with a specific key.


\textbf{Diffusion Language Models}
DLMs begin with a fixed length unknown output sequence, denoted by a set of [MASK] tokens. At each generation step, the model provides a probability distribution for its prediction for each token, and assigns a confidence score to each of its predictions. The model will unmask its predictions, and then remask the tokens with the lowest confidence scores. The model will continue this until all tokens are unmasked, producing a coherent output.


\textbf{ARLM watermarks}
Existing watermarks for ARLMs use a green-red token split. This green and red list is chosen before each word is generated, and the model's predictions are biased towards green list tokens. These watermarks often use the hash of the previous token to seed the random number generator that decides on the red-green split. This hash can not be used, since DLM generation does not always have the previous token generated.


\textbf{Existing DLM watermarks}
Gloaguen et al. \cite{gloaguen2025dlm} adapt the green-red list to diffusion models by computing the context hashes in expectation over masked tokens and applying a logit bias based on these expected hashes.
dgMARK \cite{hong2026dgmark} leaves token probabilities unchanged and instead prioritizes unmasking indices whose token prediction satisfies a keyed parity-based rule. Current methods watermark either on token distribution or unmasking order, not both.


\textbf{This work: A Joint Token-and-Order Watermark for Diffusion Language Models}
In this work, the DLM watermark is split into a token component and an order component at each diffusion generation step.
The token component alters the logit distribution at each position, and the order component modifies the confidence scores the model uses to
select which positions to unmask. We maximize the expected number of green tokens in the final output text under a specified KL per token
threshold, adding a separate bias to each component to accomplish this. The token bias is added to the logits and modeled after the energy-based bias of
Gloaguen et al. The order bias is added to each position's confidence score and equals the expected gain in green tokens
from unmasking a certain position at this particular diffusion generation step rather than at a later one. Setting either bias strength to zero effectively makes
the joint watermark a single-component one. Like ARLM watermarks, our detection algorithm counts the
number of green tokens in a section of text and performs a binomial test for significance.


\textbf{Main Contributions} Our main contributions are:
\begin{itemize}
   \item We establish a joint token-and-order watermark that applies bias to both the token distribution and the unmasking order of a DLM under a per-token KL-constrained objective, with token-only and order-only watermarks as special cases.
   \item We introduce an order component score, which measures the expected green-token gain from unmasking a position at the current generation step rather than a later one; this lets the token bias act on fixed context.
   \item We compare our joint watermark against both Gloaguen et al. and dgMARK algorithms on LLaDA-8B under the same benchmarks, green list and detector, and show higher detection at matched token-level KL.
\end{itemize}


